\doublespacing % Do not change - required

\chapter{Evaluation Plan}
\label{chEval}

%%%%%%%%%%%%%%%%%%%%%%%%%%%%%%%%%%%%%%%
% IMPORTANT
\begin{spacing}{1} %THESE FOUR
\minitoc % LINES MUST APPEAR IN
\end{spacing} % EVERY
\thesisspacing % CHAPTER
% COPY THEM IN ANY NEW CHAPTER
%%%%%%%%%%%%%%%%%%%%%%%%%%%%%%%%%%%%%%%

\section{Evaluation Objectives}

The evaluation assesses HSM-Adaptive in terms of (i) \textbf{utility} under benign training, (ii) \textbf{robustness} under Byzantine attacks, and (iii) \textbf{tolerance to heterogeneity} under non-IID client data. The goal is to demonstrate improved resilience over standard averaging without introducing excessive suppression of benign but diverse clients.

\section{Experimental Setup}
ZZ
Experiments follow a synchronous federated learning protocol with a fixed training budget (rounds, local epochs, and optimiser) shared across all methods. Standard benchmark datasets (e.g., MNIST and CIFAR-10) and consistent model architectures are used to ensure comparability. Both \textbf{IID} and \textbf{non-IID} data partitions are evaluated (e.g., label-skew partitions) to reflect realistic client heterogeneity.

\section{Baselines and Threat Model}

HSM-Adaptive is compared against FedAvg (mean aggregation) and representative robust baselines (e.g., coordinate-wise median and trimmed-mean; Krum where applicable). The threat model assumes a fraction of clients are Byzantine and can transmit arbitrary updates. Attacks include \textbf{label flipping}, \textbf{scaling}, \textbf{MSA/shuffle}, and \textbf{history-aware} variants that exploit temporal dynamics.

\section{Metrics and Protocol}

Primary metrics are \textbf{test accuracy vs.\ rounds} and \textbf{final accuracy}. Robustness diagnostics include trust/weight distributions, the fraction of near-zero weights (``floor ratio''), and effective participation
\[
\mathrm{Neff} = \frac{(\sum_i w_i)^2}{\sum_i w_i^2}.
\]
In simulated settings where attacker identities are known, detection precision/recall is reported as an auxiliary diagnostic by treating the lowest-weight clients as flagged. All experiments are repeated with fixed seeds and logged configurations to support reproducibility.

\section{Ablations and Success Criteria}

Ablation studies isolate the contribution of key components (CountSketch consensus, self-consistency, momentum alignment, and scale gating). Success is defined as: (i) minimal degradation under benign training relative to FedAvg, (ii) clear accuracy/stability improvements under attacks, and (iii) maintained Neff under non-IID settings, indicating avoidance of overly aggressive suppression.
